\documentclass[14pt]{extarticle}
% =============================================================================
% Shared Preamble for Worksheets, Homeworks, Exams & Lesson Notes
% (scripts/tex_preamble.tex)
% =============================================================================

\newif\ifdetailed
\ifdefined\setdetailed
  \setdetailed
\fi

\newif\ifwithsols
\ifdefined\setwithsols
  \setwithsols
\fi

\newif\ifwhiteblack
\ifdefined\setwhiteblack
  \setwhiteblack
\fi

% Document-class specific packages (loaded before common_preamble.tex / polyglossia)
\usepackage{fancyhdr}
\usepackage{enumitem}
\usepackage{titlesec}
\usepackage[bottom]{footmisc}
\usepackage{titling}
\usepackage{longtable} % page-breakable tables
\usepackage{float}
\usepackage[onehalfspacing]{setspace}

\ifwhiteblack
  \usepackage[a4paper, landscape, top=1.25cm, bottom=1.25cm, left=1.25cm, right=1.25cm, headheight=1.5cm, headsep=0.5cm]{geometry}
\else
  \usepackage[a4paper, top=2.4cm, bottom=2.4cm, left=2.2cm, right=2.2cm, headheight=1.4cm, headsep=0.4cm, footskip=1.0cm]{geometry}
\fi

% Shared common packages, fonts, polyglossia, callout boxes, code helpers, date & header macros
\IfFileExists{scripts/common_preamble.tex}{%
  % =============================================================================
% Shared Common Preamble (scripts/common_preamble.tex)
% Included by:
%   - scripts/tex_preamble.tex    (worksheets, homeworks, exams, lesson notes)
%   - scripts/book_preamble.tex   (two-sided Arabic book)
%   - scripts/beamer_preamble.tex (Beamer presentation slides)
% =============================================================================

\usepackage{fontspec}
\usepackage{amsmath}
\usepackage{amssymb}
\usepackage{xcolor}
\usepackage{graphicx}
\usepackage{listings}
\usepackage[most]{tcolorbox}
\usepackage{pifont}
\usepackage{minted}
\usepackage{etoolbox}
\usepackage{array}
\usepackage{booktabs}
\usepackage{tikz}
\usepackage{datetime2}

\usetikzlibrary{automata, positioning, arrows.meta, arrows, shapes.geometric, calc}

% Beamer already loads hyperref internally; load it here for article/book classes
% right before polyglossia (bidi).
\makeatletter
\@ifclassloaded{beamer}{}{%
  \usepackage{hyperref}%
  \hypersetup{hidelinks}%
  \providecommand{\LTcaptype}{table}%
}
\makeatother

% Load polyglossia (and bidi) AFTER graphics, tables, boxes, code, and hyperref
\usepackage{polyglossia}
\setmainlanguage[numerals=western]{arabic}
\setotherlanguage{english}
\newfontfamily\arabicfont[Script=Arabic]{Calibri}
\newfontfamily\arabicfontsf[Script=Arabic]{Calibri}
\newfontfamily\arabicfonttt[Script=Arabic]{Courier New}
\newfontfamily\emoji[BoldFont=*,ItalicFont=*]{Segoe UI Emoji}

% Allow slight inter-word stretch to avoid Overfull \hbox with inline \textenglish
\emergencystretch=1.5em

% -----------------------------------------------------------------------------
% Default Minted & Callout Box Definitions
% -----------------------------------------------------------------------------
\AtBeginEnvironment{minted}{\setstretch{0.95}}
\setminted{fontsize=\small, breaklines=true, vspace=2pt}

\tcbset{
  arc=2.5pt, outer arc=2.5pt,
  boxsep=1pt, left=6pt, right=6pt, top=3pt, bottom=3pt,
  toptitle=1.5pt, bottomtitle=1.5pt,
  before skip=6pt, after skip=6pt,
  fonttitle=\bfseries,
  fontupper=\setstretch{1.1}\setlength{\parskip}{2pt},
  colback=white, colframe=black, boxrule=0.6pt,
  boxcolor/.style 2 args={colback=#1, colframe=#2}
}

\newtcolorbox{boxDef}[1][]{boxcolor={blue!4!white}{blue!75!black},
  title={{\emoji📘} تعريف\ifx\\#1\\\else ~#1\fi :}}
\newtcolorbox{boxExercise}[1][]{boxcolor={cyan!4!white}{cyan!70!black},
  title={{\emoji🧩} تمرين\ifx\\#1\\\else ~#1\fi :}}
\newtcolorbox{boxExample}[1][]{boxcolor={yellow!6!white}{orange!85!black},
  title={{\emoji📝} مثال\ifx\\#1\\\else ~#1\fi :}}
\newtcolorbox{boxNote}[1][]{boxcolor={gray!7!white}{black!75},
  title={{\emoji✨} ملاحظة\ifx\\#1\\\else ~#1\fi :}}
\newtcolorbox{boxAttention}[1][]{boxcolor={magenta!5!white}{magenta!75!black},
  title={{\emoji🔔} تنبيه\ifx\\#1\\\else ~#1\fi :}}
\newtcolorbox{boxWarning}[1][]{boxcolor={red!4!white}{red!75!black},
  title={{\emoji⚡} ملاحظة هامة\ifx\\#1\\\else ~#1\fi :}}
\newtcolorbox{boxSolution}[1][]{boxcolor={green!4!white}{green!60!black},
  title={{\emoji✅} حل\ifx\\#1\\\else ~#1\fi :}}
\newtcolorbox{boxSymbol}[1][]{boxcolor={purple!4!white}{purple!70!black},
  title={{\emoji🔣} رمز\ifx\\#1\\\else ~#1\fi :}}
\newtcolorbox{boxHint}[1][]{boxcolor={teal!4!white}{teal!65!black},
  title={{\emoji💡} تلميح\ifx\\#1\\\else ~#1\fi :}}

% Compact code boxes: tighter line spacing and small vertical gaps
\tcbset{
  simplecode/.style={
    arc=2pt, outer arc=2pt,
    left=4pt, right=4pt, top=4pt, bottom=4pt,
    fontupper=\setstretch{0.95},
    before skip=6pt, after skip=6pt,
    colback=gray!5, colframe=black!50, boxrule=0.4pt
  }
}
\newenvironment{boxCode}{\begin{tcolorbox}[simplecode]}{\end{tcolorbox}}

% Code cell used inside tables (keeps code left-to-right inside RTL tables)
% Use inside a fixed-width p{...} column: \begin{pycell}\begin{minted}...\end{minted}\end{pycell}
\makeatletter
\newtcolorbox{pycell}{frame empty,colback=white,boxsep=0pt,left=2pt,right=2pt,top=2pt,bottom=2pt,width=\linewidth,before=\noindent,after=,baseline={\tcb@height-\ht\strutbox},fontupper=\setstretch{0.95}\setlength{\topsep}{0pt}\setlength{\partopsep}{0pt}\setlength{\parskip}{0pt}}
\makeatother

\newcolumntype{C}[1]{>{\centering\arraybackslash}p{#1}}

% Dotted date style (DD.MM.YYYY)
\makeatletter
\DTMnewdatestyle{dotted}{%
  \renewcommand{\DTMdisplaydate}[4]{\two@digits{##3}.\two@digits{##2}.##1}%
  \renewcommand{\DTMDisplaydate}{\DTMdisplaydate}%
}
\makeatother
\DTMsetdatestyle{dotted}

% -----------------------------------------------------------------------------
% Customizable Header / Footer Metadata
% -----------------------------------------------------------------------------
\newcommand{\headerleft}{مدرسة التسامح الشاملة}
\newcommand{\headerright}{الأستاذ محمود اغبارية}
\newcommand{\footercenter}{\thepage}

\newcommand{\setheaderleft}[1]{\renewcommand{\headerleft}{#1}}
\newcommand{\setheaderright}[1]{\renewcommand{\headerright}{#1}}
\newcommand{\setfootercenter}[1]{\renewcommand{\footercenter}{#1}}
%
}{%
  \IfFileExists{../../../scripts/common_preamble.tex}{%
    % =============================================================================
% Shared Common Preamble (scripts/common_preamble.tex)
% Included by:
%   - scripts/tex_preamble.tex    (worksheets, homeworks, exams, lesson notes)
%   - scripts/book_preamble.tex   (two-sided Arabic book)
%   - scripts/beamer_preamble.tex (Beamer presentation slides)
% =============================================================================

\usepackage{fontspec}
\usepackage{amsmath}
\usepackage{amssymb}
\usepackage{xcolor}
\usepackage{graphicx}
\usepackage{listings}
\usepackage[most]{tcolorbox}
\usepackage{pifont}
\usepackage{minted}
\usepackage{etoolbox}
\usepackage{array}
\usepackage{booktabs}
\usepackage{tikz}
\usepackage{datetime2}

\usetikzlibrary{automata, positioning, arrows.meta, arrows, shapes.geometric, calc}

% Beamer already loads hyperref internally; load it here for article/book classes
% right before polyglossia (bidi).
\makeatletter
\@ifclassloaded{beamer}{}{%
  \usepackage{hyperref}%
  \hypersetup{hidelinks}%
  \providecommand{\LTcaptype}{table}%
}
\makeatother

% Load polyglossia (and bidi) AFTER graphics, tables, boxes, code, and hyperref
\usepackage{polyglossia}
\setmainlanguage[numerals=western]{arabic}
\setotherlanguage{english}
\newfontfamily\arabicfont[Script=Arabic]{Calibri}
\newfontfamily\arabicfontsf[Script=Arabic]{Calibri}
\newfontfamily\arabicfonttt[Script=Arabic]{Courier New}
\newfontfamily\emoji[BoldFont=*,ItalicFont=*]{Segoe UI Emoji}

% Allow slight inter-word stretch to avoid Overfull \hbox with inline \textenglish
\emergencystretch=1.5em

% -----------------------------------------------------------------------------
% Default Minted & Callout Box Definitions
% -----------------------------------------------------------------------------
\AtBeginEnvironment{minted}{\setstretch{0.95}}
\setminted{fontsize=\small, breaklines=true, vspace=2pt}

\tcbset{
  arc=2.5pt, outer arc=2.5pt,
  boxsep=1pt, left=6pt, right=6pt, top=3pt, bottom=3pt,
  toptitle=1.5pt, bottomtitle=1.5pt,
  before skip=6pt, after skip=6pt,
  fonttitle=\bfseries,
  fontupper=\setstretch{1.1}\setlength{\parskip}{2pt},
  colback=white, colframe=black, boxrule=0.6pt,
  boxcolor/.style 2 args={colback=#1, colframe=#2}
}

\newtcolorbox{boxDef}[1][]{boxcolor={blue!4!white}{blue!75!black},
  title={{\emoji📘} تعريف\ifx\\#1\\\else ~#1\fi :}}
\newtcolorbox{boxExercise}[1][]{boxcolor={cyan!4!white}{cyan!70!black},
  title={{\emoji🧩} تمرين\ifx\\#1\\\else ~#1\fi :}}
\newtcolorbox{boxExample}[1][]{boxcolor={yellow!6!white}{orange!85!black},
  title={{\emoji📝} مثال\ifx\\#1\\\else ~#1\fi :}}
\newtcolorbox{boxNote}[1][]{boxcolor={gray!7!white}{black!75},
  title={{\emoji✨} ملاحظة\ifx\\#1\\\else ~#1\fi :}}
\newtcolorbox{boxAttention}[1][]{boxcolor={magenta!5!white}{magenta!75!black},
  title={{\emoji🔔} تنبيه\ifx\\#1\\\else ~#1\fi :}}
\newtcolorbox{boxWarning}[1][]{boxcolor={red!4!white}{red!75!black},
  title={{\emoji⚡} ملاحظة هامة\ifx\\#1\\\else ~#1\fi :}}
\newtcolorbox{boxSolution}[1][]{boxcolor={green!4!white}{green!60!black},
  title={{\emoji✅} حل\ifx\\#1\\\else ~#1\fi :}}
\newtcolorbox{boxSymbol}[1][]{boxcolor={purple!4!white}{purple!70!black},
  title={{\emoji🔣} رمز\ifx\\#1\\\else ~#1\fi :}}
\newtcolorbox{boxHint}[1][]{boxcolor={teal!4!white}{teal!65!black},
  title={{\emoji💡} تلميح\ifx\\#1\\\else ~#1\fi :}}

% Compact code boxes: tighter line spacing and small vertical gaps
\tcbset{
  simplecode/.style={
    arc=2pt, outer arc=2pt,
    left=4pt, right=4pt, top=4pt, bottom=4pt,
    fontupper=\setstretch{0.95},
    before skip=6pt, after skip=6pt,
    colback=gray!5, colframe=black!50, boxrule=0.4pt
  }
}
\newenvironment{boxCode}{\begin{tcolorbox}[simplecode]}{\end{tcolorbox}}

% Code cell used inside tables (keeps code left-to-right inside RTL tables)
% Use inside a fixed-width p{...} column: \begin{pycell}\begin{minted}...\end{minted}\end{pycell}
\makeatletter
\newtcolorbox{pycell}{frame empty,colback=white,boxsep=0pt,left=2pt,right=2pt,top=2pt,bottom=2pt,width=\linewidth,before=\noindent,after=,baseline={\tcb@height-\ht\strutbox},fontupper=\setstretch{0.95}\setlength{\topsep}{0pt}\setlength{\partopsep}{0pt}\setlength{\parskip}{0pt}}
\makeatother

\newcolumntype{C}[1]{>{\centering\arraybackslash}p{#1}}

% Dotted date style (DD.MM.YYYY)
\makeatletter
\DTMnewdatestyle{dotted}{%
  \renewcommand{\DTMdisplaydate}[4]{\two@digits{##3}.\two@digits{##2}.##1}%
  \renewcommand{\DTMDisplaydate}{\DTMdisplaydate}%
}
\makeatother
\DTMsetdatestyle{dotted}

% -----------------------------------------------------------------------------
% Customizable Header / Footer Metadata
% -----------------------------------------------------------------------------
\newcommand{\headerleft}{مدرسة التسامح الشاملة}
\newcommand{\headerright}{الأستاذ محمود اغبارية}
\newcommand{\footercenter}{\thepage}

\newcommand{\setheaderleft}[1]{\renewcommand{\headerleft}{#1}}
\newcommand{\setheaderright}[1]{\renewcommand{\headerright}{#1}}
\newcommand{\setfootercenter}[1]{\renewcommand{\footercenter}{#1}}
%
  }{%
    \input{../../scripts/common_preamble.tex}%
  }%
}

\pagestyle{fancy}

% Prevent widow/orphan lines
\widowpenalty=10000
\clubpenalty=10000

% Compact section headings and default list spacing so pages do not waste vertical space
\titlespacing*{\section}{0pt}{1.4ex plus 0.4ex minus 0.2ex}{0.7ex plus 0.2ex}
\titlespacing*{\subsection}{0pt}{1.1ex plus 0.3ex minus 0.2ex}{0.5ex plus 0.1ex}
\titlespacing*{\subsubsection}{0pt}{0.9ex plus 0.2ex minus 0.1ex}{0.4ex plus 0.1ex}
\setlist{topsep=4pt, partopsep=0pt, parsep=2pt, itemsep=4pt}

\gappto\captionsarabic{\renewcommand{\contentsname}{جدول المحتويات}}

% Legacy listings settings
\lstset{
  language=[Sharp]C,
  numbers=left,
  stepnumber=1,
  numbersep=8pt,
  frame=single,
  basicstyle=\ttfamily\small
}
\ifwhiteblack
  \lstset{
    keywordstyle=\bfseries,
    stringstyle=\ttfamily,
    commentstyle=\itshape
  }
\else
  \lstset{
    keywordstyle=\color{blue},
    stringstyle=\color{red},
    commentstyle=\color{green!50!black}
  }
\fi

% White-and-black landscape presentation mode overrides
\ifwhiteblack
  \AtBeginDocument{
    \LARGE
    \usemintedstyle{bw}
  }
  \setminted{vspace=2pt}
  \tcbset{
    colback=white, colframe=black, boxrule=1pt,
    boxcolor/.style 2 args={colback=white, colframe=black},
    simplecode/.append style={colback=white, colframe=black, boxrule=1pt}
  }
\fi

% Default to empty author to avoid "LaTeX Warning: No \author given" when omitted
\author{}

% Compact title block
\makeatletter
\renewcommand{\@maketitle}{%
  \begin{center}%
    {\ifwhiteblack\Huge\else\LARGE\fi \bfseries \@title \par}%
    \ifdefempty{\@author}{}{%
      \vskip 0.15em%
      {\large \@author \par}%
    }%
    {\small \hfill \textenglish{\DTMtoday} \par}%
  \end{center}%
  \vspace{-0.3cm}%
}
\makeatother

\apptocmd{\maketitle}{\thispagestyle{fancy}}{}{}

% Apply page headers and footer (using \headerleft, \headerright, \footercenter from common_preamble.tex)
\newcommand{\applyfancyheaders}{%
  \fancyhf{}%
  \fancyhead[L]{\textarabic{\headerleft}}%
  \fancyhead[R]{\textarabic{\headerright}}%
  \fancyfoot[C]{\textarabic{\footercenter}}%
}

\fancypagestyle{plain}{\applyfancyheaders}
\applyfancyheaders

\setcounter{tocdepth}{3} % only section, subsection, and subsubsection in TOC

\newcommand{\importpdfpage}[2]{%
    \newgeometry{left=0cm,right=0cm,top=0cm,bottom=0cm}
    \thispagestyle{empty}
    \includegraphics[page=#2,width=0.95\paperwidth,height=0.95\paperheight,keepaspectratio]{#1}
    \restoregeometry
}

\newcommand{\insertFullImg}[1]{%
    \noindent
    \makebox[\textwidth][c]{
        \includegraphics[width=0.9\paperwidth,keepaspectratio]{#1}%
    }%
}%


\title{مشروع شركة "السعادة للجميع"}

\begin{document}

\maketitle
\thispagestyle{fancy}

% ==============================================================
\section*{المرحلة الثالثة: التحويل لأعلى وإعادة تعريف العمليات (\textenglish{virtual} و \textenglish{override})}
% ==============================================================

\textbf{قواعد حساب الرواتب ونقاط التصويت واستحقاق المكافأة} \\
تُتخذ القرارات في الشركة عبر التصويت باستخدام \textbf{نقاط التأثير}، كما يتحدد الراتب الشهري والمكافأة السنوية لكل موظف بحسب نمطه وأقدميته كما يلي:

{\renewcommand{\arraystretch}{1.4}
\small
\begin{longtable}{|p{0.20\textwidth}|p{0.37\textwidth}|p{0.35\textwidth}|}
\hline
\textbf{نمط الموظف} & \textbf{الراتب الشهري (\textenglish{GetSalary})} & \textbf{نقاط التصويت (\textenglish{GetScore})} \\
\hline
\textbf{موظف عادي} \newline (\textenglish{Employee}) &
راتب أساسي $5000$ ش.ج + $500$ ش.ج عن كل سنة أقدمية: \newline
$5000 + 500 \cdot \text{seniority}$ &
$4$ نقاط أساسية + نقطة عن كل سنة أقدمية: \newline
$4 + \text{seniority}$ \\
\hline
\textbf{تقني حواسيب} \newline (\textenglish{Technician}) &
راتب الموظف العادي + $250$ ش.ج عن كل $10$ حواسيب كاملة: \newline
$\text{base} + \lfloor \frac{\text{computers}}{10} \rfloor \cdot 250$ &
نقاط الموظف العادي + نقطة عن كل حاسوب: \newline
$\text{base} + \text{computers}$ \\
\hline
\textbf{عضو لجنة عمال} \newline (\textenglish{UnionMember}) &
يماثل راتب الموظف العادي تماماً. &
\textbf{ضعف} نقاط الموظف العادي + نقطتان عن كل سنة باللجنة: \newline
$2 \cdot \text{base} + 2 \cdot \text{unionYears}$ \\
\hline
\textbf{مسؤول مشاريع} \newline (\textenglish{Supervisor}) &
راتب الموظف العادي + $100$ ش.ج عن كل موظف فعلي في فريقه: \newline
$\text{base} + 100 \cdot \text{teamCount}$ &
نقاط الموظف العادي + مجموع نقاط كافة أعضاء فريقه الفعليين. \\
\hline
\end{longtable}}

\filbreak

\begin{enumerate}[label=\textbf{تمرين \arabic*:}, wide=0pt, itemsep=2.2em]

    \item \textbf{تطبيق العمليات القابلة لإعادة التعريف (\textenglish{GetSalary()} و \textenglish{GetScore()}) واستدعاء \textenglish{base}} \\
    أضف إلى الفئة الأساسية \textenglish{Employee} العمليتين العامتين القابلتين لإعادة التعريف:
    \begin{itemize}
        \item \textenglish{public virtual double GetSalary()}
        \item \textenglish{public virtual int GetScore()}
    \end{itemize}
    وأعد تعريفهما (\textenglish{override}) في الفئات الفرعية \textenglish{Technician} و \textenglish{UnionMember} و \textenglish{Supervisor} وفق الجدول أعلاه، مع الحرص على استدعاء عملية فئة الأب عبر \textenglish{base.GetSalary()} و \\ \textenglish{base.GetScore()} لتجنب تكرار المعادلات الحسابية. \\
    \textbf{الفحص في \textenglish{Main} (التحويل لأعلى):} \\
    أنشئ في \textenglish{Main} مؤشرات من نمط الأب \textenglish{Employee} تشير إلى كائنات من الفئات الوارثة (مثلاً: \textenglish{Employee e1 = new Technician(...);} و \textenglish{Employee e2 = new UnionMember(...);} و \textenglish{Employee e3 = new Supervisor(...);})، ثم استدعِ عليها \textenglish{GetSalary()} و \textenglish{GetScore()} و \textenglish{ToString()} واطبع النتائج.
    \filbreak

    \item \textbf{إضافة عملية استحقاق المكافأة السنوية (\textenglish{YearlyBonus})} \\
    أضف لمشروع الشركة عملية عامة قابلة لإعادة التعريف تعريفها: \\
    \textenglish{public virtual bool YearlyBonus()}
    \begin{itemize}
        \item تُعرّف العملية لأول مرة في الفئة \textenglish{Employee} وتعيد \textenglish{true} إذا كانت أقدمية الموظف $5$ سنوات على الأقل ($\text{seniority} \ge 5$)، وإلا تعيد \textenglish{false}.
        \item في الفئة \textenglish{Technician} يُعاد تعريف العملية (\textenglish{override}) وتعيد \textenglish{true} إذا كان التقني يحقق شرط الأقدمية (عبر استدعاء \textenglish{base.YearlyBonus()}) \textbf{وأيضًا} يتولى مسؤولية $100$ حاسوب على الأقل ($\text{computers} \ge 100$)، وإلا تعيد \textenglish{false}.
        \item في الفئتين \textenglish{UnionMember} و \textenglish{Supervisor} يُعاد تعريف العملية (\textenglish{override}) بحيث تعيد \textenglish{false} دائماً (لأن أعضاء اللجنة ومسؤولي المشاريع غير مشمولين في هذه المكافأة).
    \end{itemize}
    \textbf{الفحص في \textenglish{Main}:} \\
    استدعِ في \textenglish{Main} العملية \textenglish{YearlyBonus()} عبر مؤشرات من نمط \textenglish{Employee} تشير إلى كائنات مختلفة واطبع النتائج.
    \filbreak

    \item \textbf{بناء فئة إدارة الشركة (\textenglish{Company}) وإضافة موظفين (\textenglish{AddEmployee})} \\
    أضف فئة جديدة لمشروعنا باسم \textenglish{Company}. تمثّل هذه الفئة مجموعة موظفي الشركة من جميع الأنماط، بالاعتماد على حفظ مؤشرات لكائنات الأبناء داخل مصفوفة من نمط الأب \textenglish{Employee[]}:

    {\renewcommand{\arraystretch}{1.2}
    \small
    \begin{longtable}{|>{\raggedleft\arraybackslash}p{0.4\textwidth}|>{\raggedright\arraybackslash}p{0.5\textwidth}|}
    \hline
    \multicolumn{2}{|c|}{\textbf{الفئة: \textenglish{Company}}} \\
    \hline
    صفات: ثابت يمثل الحد الأقصى للموظفين في الشركة، مصفوفة لحفظ موظفي الشركة، وعدد الموظفين الفعليين حالياً فيها. & \textenglish{public const int MAX = 100;} \newline \textenglish{private Employee[] employees;} \newline \textenglish{private int current;} \\
    \hline
    عملية بنائية تنشئ المصفوفة بحجم \textenglish{MAX} وتهيئ \textenglish{current} إلى $0$. & \textenglish{public Company()} \\
    \hline
    عملية تعيد عدد الموظفين الحاليين في الشركة. & \textenglish{GetCurrent()} \\
    \hline
    عملية تتلقى موظفاً من أي نمط وتضيفه إلى مصفوفة موظفي الشركة: تفحص أولاً إن كان هناك متسع في المصفوفة؛ إذا نعم، تضيف الموظف وتعيد \textenglish{true}، وإلا فلا تضيفه وتعيد \textenglish{false}. & \textenglish{public bool AddEmployee(Employee e)} \\
    \hline
    \end{longtable}}

    \vspace{1em}
    \noindent \textbf{الفحص في \textenglish{Main}:} \\
    أنشئ كائناً من الفئة \textenglish{Company} في \textenglish{Main}، وأضف إليه عدة موظفين من الأنماط الأربعة (\textenglish{Employee} و \textenglish{Technician} و \textenglish{UnionMember} و \textenglish{Supervisor})، واطبع عدد الموظفين الحاليين.
    \filbreak

    \item \textbf{عمليات تعتمد على \textenglish{virtual} و \textenglish{override} في الفئة \textenglish{Company}}
    \begin{enumerate}[label=\textbf{(\arabic*)}]
        \item \textbf{البحث عن بيانات موظف حسب رقم الهوية (\textenglish{GetEmployeeDetails}):} \\
        اكتب في الفئة \textenglish{Company} عملية باسم \textenglish{public string GetEmployeeDetails(string id)} تتلقى رقم هوية موظف (افترض أنه فريد) وتعيد سلسلة نصية تصف بيانات هذا الموظف (باستدعاء \textenglish{ToString()}). إذا لم يوجد موظف بهذا الرقم في الشركة، تعيد العملية نصاً فارغاً \textenglish{""}.
        \item \textbf{مجموع الرواتب الشهرية (\textenglish{TotalSalary}):} \\
        اكتب في الفئة \textenglish{Company} عملية باسم \textenglish{public double TotalSalary()} تحسب وتعيد مجموع الرواتب الشهرية لجميع موظفي الشركة (باستدعاء \textenglish{GetSalary()} لكل موظف في المصفوفة).
        \item \textbf{الموظفون المستحقون للمكافأة (\textenglish{GetBonusEmployees}):} \\
        اكتب في الفئة \textenglish{Company} عملية باسم \textenglish{public Employee[] GetBonusEmployees()} تعيد مصفوفة (\textenglish{Employee[]}) تضم جميع موظفي الشركة المستحقين للمكافأة السنوية (بحسب نتيجة العملية \textenglish{YearlyBonus()}).
    \end{enumerate}
    \textbf{الفحص في \textenglish{Main}:} \\
    في \textenglish{Main}، استدعِ \textenglish{GetEmployeeDetails} لرقم هوية موجود وآخر غير موجود، ثم اطبع مجموع الرواتب الكلي عبر \textenglish{TotalSalary()}، واستدعِ \textenglish{GetBonusEmployees()} واطبع بيانات الموظفين المستحقين للمكافأة.
    \filbreak

\end{enumerate}

\vspace{1cm}
\begin{flushleft}
أرجو لكم وقتًا ممتعًا.

الأستاذ محمود اغبارية.
\end{flushleft}

\end{document}
